% Generated draft: requires mathematical and source review.
\documentclass[11pt]{article}
\usepackage[margin=25mm]{geometry}
\usepackage{fontspec}
\setmainfont{Latin Modern Roman}
\usepackage{amsmath,amssymb,mathtools,unicode-math}
\setmathfont{Latin Modern Math}
\usepackage{enumitem,needspace,xurl,hyperref}
\setlength{\parindent}{0pt}
\setlength{\parskip}{5pt}
\setlength{\emergencystretch}{3em}
\newcommand{\N}{\mathbb{N}}
\newcommand{\Z}{\mathbb{Z}}
\newcommand{\Q}{\mathbb{Q}}
\newcommand{\R}{\mathbb{R}}
\newcommand{\C}{\mathbb{C}}
\newcommand{\F}{\mathbb{F}}
\newcommand{\A}{\mathbb{A}}
\newcommand{\PP}{\mathbb{P}}
\newcommand{\E}{\mathbb{E}}
\newcommand{\eps}{\varepsilon}
\newcommand{\dd}{\,\mathrm{d}}
\begin{document}
\section*{UnsolvedMath ID: 30004318}
\textbf{Source problem code:} catalogue entry
\subsection{Definitions and mathematical statement}
The supplied question is: ``Is the known list of rank-one groups, consisting of examples of algebraic origin such as groups related to algebraic or classical groups and small-characteristic variants, complete?''

The record supplies neither a definition of a rank-one group nor the list whose completeness is questioned. It also does not specify the permitted fields, characteristics, additional hypotheses, or notion of equivalence. Consequently, a precise classification statement cannot be recovered from this record alone.

A conditional formulation is as follows. Let \(\mathcal C\) denote the intended class of rank-one groups, let \(\mathcal L\subseteq\mathcal C\) denote the intended known list, and let \(\sim\) denote the intended equivalence relation on \(\mathcal C\). Once these data are specified, the completeness question is whether
\[
\forall G\in\mathcal C\;\exists H\in\mathcal L\quad G\sim H.
\]
Here \(\mathcal C\), \(\mathcal L\), and \(\sim\) are unresolved placeholders, not definitions recovered from the source. In particular, the record does not determine whether algebraic origin restricts the groups to be classified or merely describes the currently known examples. The reference to small-characteristic variants supplies no precise characteristic restrictions or constructions.
\subsection{Short English statement}
The question asks whether every rank-one group under the intended hypotheses occurs in the intended known list. The definition of rank one, the list, the hypotheses, and the equivalence relation must be recovered before this becomes a precise mathematical statement. The supplied record labels the problem open; that label is not independently verified.
\subsection{Sources}
\begin{itemize}
\item Local UnsolvedMath registry, \texttt{data/unsolvedmath/problems.json}, permanent ID \(30004318\), title ``Classification of Rank-One Groups,'' source problem code \texttt{catalogue entry}. The supplied record belongs to the Oberwolfach Reports Open Problems collection. No individual proposer is supplied.
\item Supplied report citation: \emph{Mini-Workshop: Rank One Groups and Exceptional Algebraic Groups} (2020), Oberwolfach Reports, pp.\ 3243--3269, DOI \texttt{10.4171/owr/2019/52}. \url{https://doi.org/10.4171/owr/2019/52}. The background describes these reports as documenting mathematical workshops and their open-problem discussions.
\item Supplied literature source: ``Oberwolfach Reports, rank-one groups.'' \url{https://ems.press/content/serial-article-files/46832}.
\item Literature assessment dated exactly \texttt{2026-08-22}: the supplied curation labels the problem open and describes completeness of the known list as an explicit open question. Its description of the classification as incomplete does not supply a counterexample to completeness. This is supplied dated curation derived from the OWR material, not an independently verified literature assessment; the record itself requests independent verification of time-sensitive status and cited claims.
\end{itemize}
\end{document}
